\subsection{Convergence for all obstacles}
\label{subsec:curve-all-obstacle-topology}

We keep the compact curve $Y$ and its positive form $\theta$ of area $A$. The topology $\tau^{\mathrm{all}}_\theta$ of \cref{subsec:all-obstacle-topology} is generated by the topologies $\tau_{\theta,F}$ for all smooth obstacles $F$, in which the ordinary and ideal envelopes below $F$ replace those below zero. On a curve the envelopes below $F$ are still determined by the polar measure, and their ideal counterparts by its atomic part. This gives a description of $\tau^{\mathrm{all}}_\theta$ by a single pseudometric and allows a comparison with polar total variation.

Choose, as in \cref{prop:all-obstacle-uniformity}, a countable family $\{F_k\}_{k\geq1}\subset C^\infty(Y,\mathbb R)$ that is dense for the uniform norm, with $F_1=0$. Write $d_{O,F}$ and $d_{I,F}$ for the $L^1$ distances between the ordinary envelopes $P_{\theta,F}[\cdot]$ and between the ideal envelopes $P_{\theta,F}[\cdot]_{\mathcal I}$, which define $\tau_{\theta,F}$ (\cref{sec:ordinary-type-topology}), and put
\begin{equation}\label{eq:curve-all-obstacle-distance}
 \begin{split}
 d^{\mathrm{all}}_\theta([u],[v])
 &=\|D_u-D_v\|_{\mathrm{TV}}\\
 &\quad+\sum_{k=1}^\infty2^{-k}\bigl(\min\{1,d_{O,F_k}([u],[v])\}\\
 &\qquad\qquad+\min\{1,d_{I,F_k}([u],[v])\}\bigr).
 \end{split}
\end{equation}

\begin{proposition}\label{prop:curve-all-obstacle-topology}\index{Singularity type!all-obstacle curve pseudometric}
The pseudometric $d^{\mathrm{all}}_\theta$ induces the uniformity of $\tau^{\mathrm{all}}_\theta$. Two classes have distance zero if and only if they have the same polar curvature. The closure of the analytic classes is $\{[u]:\mu_u=D_u\}=\mathcal G_\theta$, where $\mathcal G_\theta=\{[u]:P_\theta[u]=P_\theta[u]_{\mathcal I}\}$ as in \eqref{eq:ordinary-topology-analytic-closure}. On each set $\mathcal C_s$, $0\leq s\leq A$, the pseudometric $d^{\mathrm{all}}_\theta$ is a complete metric inducing the topology of total variation of the atomic measures. On ordinary models, polar total variation implies convergence for $d^{\mathrm{all}}_\theta$.
\end{proposition}
\begin{proof}
\emph{Step 1: envelopes below an obstacle.} For any smooth $F$, translating the competitors by $F$ gives
\[
 P_{\theta,F}[u]-F=P_{\theta+\ddc F}[u-F],
\]
and likewise for the ideal envelope. The background $\theta+\ddc F$ on the right is a smooth real form of area $A$, though it need not be positive. We claim that these two envelopes are the models of $\mu_u$ and of $D_u$, respectively, in that background, in the sense of \cref{prop:curve-arbitrary-background-models} in \cref{sec:curve-model-backgrounds}; the proof of that proposition does not use the present subsection. Set $\mu=\mu_u$ and let $w_\mu$ be its Green potential, and put $r=P_{\theta+\ddc F}[w_\mu-F]$. By the gauge bounds of \cref{prop:i-envelope-gauge}, $r+F$ differs from $P_\theta[w_\mu]=p_\mu\sim w_\mu$ by a bounded function, so $r\sim w_\mu-F$. Hence $r=P_{\theta+\ddc F}[r]$ is an ordinary model, and its polar measure is $\mu$. Moreover $u-F\preceq r$, and both have mass $A-\mu(Y)$. If this mass is positive, \cref{cor:full-mass-envelope-characterization}, applied in the big background $\theta+\ddc F$, shows that $P_{\theta+\ddc F}[u-F]=r$. If it is zero, the curvature measures of $u-F$ and $w_\mu-F$ both equal $\mu$, so these functions differ by a constant, and again $P_{\theta+\ddc F}[u-F]=r$. The assertion for the ideal envelope follows in the same way from ideal projection and the classification of atomic models. It also follows that the envelope $P_{\theta,F}[u]$ is the upper envelope of all $v\leq F$ with $\theta+\ddc v\geq\mu$.

\emph{Step 2: the uniformity.} By \cref{prop:curve-arbitrary-background-models}, polar total variation gives $L^1$ convergence of these models for every fixed $F$. The same assertion for atomic measures gives continuity of every ideal-envelope coordinate with respect to the total variation of the atomic parts. The controlled ideal relations are also implied, uniformly, by atomic total variation: for fixed $t>0$ and $a>1$, the bound $\|D-E\|_{\mathrm{TV}}<(a-1)/(at)$ guarantees both inclusions in \eqref{eq:controlled-ideal-entourage}, by the formula for multiplier ideals on a curve. Indeed, if $tE(\{x\})\geq1$, the bound gives $aD(\{x\})>E(\{x\})$; otherwise the target ideal is the unit ideal. Interchange $D$ and $E$ for the other inclusion. Each divisorial coordinate is the atomic total variation. Therefore the envelope terms in the series, together with the atomic term, generate exactly the uniformity of $\tau^{\mathrm{all}}_\theta$: finitely many coordinates control the sum up to its uniformly small tail, and the sum controls each fixed coordinate.

\emph{Step 3: zero distance and analytic closure.} By Step~1 with $F=0$, distance zero forces equal polar measures; the classification at every obstacle proves the converse. The same classification shows that the ordinary and ideal envelopes agree at every $F$ if and only if $\mu_u=D_u$. Hence \cref{thm:all-obstacle-analytic-closure} identifies the closure of the analytic classes.

\emph{Step 4: the sets $\mathcal C_s$.} For a $d^{\mathrm{all}}_\theta$-Cauchy sequence in $\mathcal C_s$, the atomic measures converge in total variation to a measure of mass at most $s$, which is again atomic. The convergence of models in Step~2 gives convergence of every obstacle coordinate, and the summable tail then gives convergence for $d^{\mathrm{all}}_\theta$. Conversely, convergence for $d^{\mathrm{all}}_\theta$ implies convergence of the atomic parts, because of the first summand. This proves completeness and the equality of topologies on $\mathcal C_s$. Finally, polar total variation controls the atomic part and all ordinary-envelope coordinates, which proves the last assertion.
\end{proof}

On a curve, membership in the relative envelope class of the ideal envelope has the same polar criterion, including for ideal envelopes of zero mass.

\Needspace{18\baselineskip}
\begin{corollary}\label{cor:curve-relative-envelope-criterion}\index{Relative envelope class!polar criterion on a curve}\index{Relative envelope class!zero-mass curve model}
For $u\in\PSH(Y,\theta)$ put $q=P_\theta[u]_{\mathcal I}$. The following are equivalent:
\[
 \begin{gathered}
 u\in\mathcal E_{\mathrm{env}}(Y,\theta,q),\qquad \mu_u=D_u,\\
 P_\theta[u]=P_\theta[u]_{\mathcal I},\\
 P_{\theta,F}[u]=P_{\theta,F}[u]_{\mathcal I}
       \quad\text{for every smooth real }F.
 \end{gathered}
\]
\end{corollary}
\begin{proof}
The classifications of ordinary and ideal models in \cref{lem:curve-polar-model,lem:curve-atomic-ideal-model} show that $P_\theta[u]=P_\theta[u]_{\mathcal I}$ if and only if $\mu_u=D_u$. Step~1 of the proof of \cref{prop:curve-all-obstacle-topology} shows that this is equivalent to equality of the envelopes at every smooth obstacle. Membership of $u$ in $\mathcal E_{\mathrm{env}}(Y,\theta,q)$ implies the latter by \cref{prop:ideal-obstacle-reprojection}.

Conversely, assume that $\mu_u=D_u$. By \cref{lem:curve-atomic-ideal-model}, $\mu_q=D_u$, so both singularity volumes of $q$ equal $A-D_u(Y)$. If $D_u(Y)<A$, then $P_\theta[u]=P_\theta[u]_{\mathcal I}$, and \cref{prop:conditional-i-characterization} puts $u$ in $\mathcal E_{\mathrm{env}}(Y,\theta,q)$.

If $D_u(Y)=A$, the positive curvature measures $T_u$ and $T_q$ both dominate $D_u$ and both have total mass $A$. Hence $T_u=T_q=D_u$. The distribution $u-q$ is harmonic, and therefore constant on the compact connected curve. Thus $u=q+C$ for a constant $C$, and \cref{prop:relative-elementary} gives $u\in\mathcal E_{\mathrm{env}}(Y,\theta,q)$ also in this case.
\end{proof}

\begin{proposition}\label{prop:curve-all-obstacle-incomplete}
For the Fubini--Study form $\theta$ of area one on $\mathbb P^1$, the pseudometric $d^{\mathrm{all}}_\theta$ is incomplete on $\mathfrak S(\mathbb P^1,\theta)$: some Cauchy sequence of ordinary models of one fixed positive mass with trivial multiplier ideals has no limit. On ordinary models, its topology is strictly weaker than that of polar total variation and strictly stronger than $\tau_{\theta,0}$.
\end{proposition}
\begin{proof}
Use the diffuse polar probability measures $\sigma_j$ concentrating at a point $x$ and the models $p_j=p_{c\sigma_j}$ from the proof of \cref{prop:curve-ordinary-topology-incomplete}, with $0<c<1$. Fix a smooth obstacle $F$ and put $q_{y,F}=P_{\theta,F}[q_{c\delta_y}]$. Rotations preserving $\theta$ and sending $x$ to $y$ can be chosen to depend continuously on $y$ near $x$; they transport the obstacle $F$ to its composition with the rotation. Since envelopes below two obstacles differ by at most the uniform distance between the obstacles, continuity of rotations in $L^1$ gives $q_{y,F}\to q_{x,F}$ in $L^1$ as $y\to x$.

The uniform $L^1$ bound near $x$ makes the average $w_{j,F}=\int q_{y,F}\,\mathrm d\sigma_j(y)$ well defined. It lies below $F$, its curvature current dominates $c\sigma_j$, and it converges to $q_{x,F}$ in $L^1$. The description of envelopes below obstacles in Step~1 of the proof of \cref{prop:curve-all-obstacle-topology} gives $w_{j,F}\leq P_{\theta,F}[p_j]\leq F$. These bounds and normalized compactness give $L^1$ precompactness. Every cluster point of the envelopes has curvature at least $c\delta_x$, so it is at most $q_{x,F}$, and the lower bound gives the reverse inequality. Hence $P_{\theta,F}[p_j]\to q_{x,F}$ for every $F$. The ideal envelopes of the $p_j$ below $F$ all equal $P_\theta(F)$, since their multiplier ideals are trivial, and their atomic parts vanish. Thus $[p_j]$ is a $d^{\mathrm{all}}_\theta$-Cauchy sequence. Since $F_1=0$, a limit would also be a $d^{\mathrm{ot}}_\theta$-limit, which \cref{prop:curve-ordinary-topology-incomplete} excludes.

For the strict comparison with polar total variation, use the rotated diffuse models $\widetilde p_j=p\circ R_j^{-1}$ from the end of \cref{subsec:curve-ordinary-topology}. For every smooth $F$,
\[
 P_{\theta,F}[\widetilde p_j]
   =P_{\theta,F\circ R_j}[p]\circ R_j^{-1}.
\]
The same obstacle bound and the $L^1$ continuity of rotations give convergence to $P_{\theta,F}[p]$. All atomic parts vanish, so $d^{\mathrm{all}}_\theta([\widetilde p_j],[p])\to0$, whereas the polar total-variation distance remains $2c$. \Cref{prop:curve-all-obstacle-topology} gives the implication in the other direction.

Finally, the sequence in \cref{prop:ordinary-topology-gauge-failure} converges for the zero obstacle but not for some smooth obstacle. By Step~1 of the proof of \cref{prop:curve-all-obstacle-topology}, replacing its members by their ordinary models does not change any obstacle envelope, so $d^{\mathrm{all}}_\theta$ is strictly stronger than $\tau_{\theta,0}$ already on ordinary models.
\end{proof}
